\documentclass[11pt, a4paper]{article}

\usepackage[T1]{fontenc}
\usepackage[utf8]{inputenc}
\usepackage{tgpagella}
\usepackage{microtype}
\usepackage[spanish, es-noshorthands]{babel}

\usepackage{amsmath, amssymb, amsfonts}
\usepackage[table, xcdraw]{xcolor}
\usepackage{booktabs}
\usepackage[most]{tcolorbox}
\usepackage[margin=2.5cm]{geometry}
\usepackage{fancyhdr}

\pagestyle{fancy}
\fancyhf{}
\setlength{\headheight}{16pt}
\lhead{\textbf{IMT342}}
\chead{Ejercicios: Singularidades}
\rhead{Semana 9 - Sesión 1}
\lfoot{Prof: M.Sc. Ing. Bernardo Quiroga Turdera}
\rfoot{Página \thepage}
\renewcommand{\headrulewidth}{0.4pt}

\definecolor{ucbblue}{RGB}{0, 51, 102}

\begin{document}

\begin{center}
    \Large\textbf{\textcolor{ucbblue}{Hoja de Ejercicios: Rango y Singularidad Cinemática}}\\[0.2cm]
\end{center}

\vspace{0.4cm}

\section*{Ejercicio 1 — Mapeo de Velocidad y Significado[cite: 13]}
Dado el Jacobiano evaluado y el vector de velocidades articulares:
\begin{equation*}
    J = \begin{bmatrix} 0 & -1 \\ 2 & 1 \end{bmatrix}, \qquad \dot{\mathbf{q}} = \begin{bmatrix} 1 \\ 0 \end{bmatrix} \text{[cite: 13]}
\end{equation*}
\textbf{a)} Calcule la velocidad del efector final $\dot{\mathbf{x}} = J\dot{\mathbf{q}}$. \\
\textbf{b)} Explique físicamente qué representa la primera columna de la matriz $J$.
\vspace{2.5cm}

\section*{Ejercicio 2 — Clasificación de Rango[cite: 13]}
Dadas las siguientes matrices:
\begin{equation*}
    A = \begin{bmatrix} 1 & 0 \\ 0 & 1 \end{bmatrix}, \qquad B = \begin{bmatrix} 1 & -2 \\ -2 & 4 \end{bmatrix} \text{[cite: 13]}
\end{equation*}
Calcule analíticamente el rango de ambas matrices e indique si representan un estado regular o singular.
\vspace{3cm}

\section*{Ejercicio 3 — Criterio General vs Cuadrado[cite: 13]}
Escriba la definición universal matemática de una singularidad cinemática mediante el rango, y explique por qué la ecuación $\det(J) = 0$ es un caso especial y no la regla general.
\vspace{3cm}

\section*{Ejercicio 4 — Singularidad Geométrica 2R[cite: 13]}
Para un manipulador 2R con $a_1=a_2=1$ y $q_1=0$, compare las posturas correspondientes a $q_2 = \pi/2$ y $q_2 = 0$.
Para la configuración singular, identifique:
\begin{enumerate}
    \item La relación lineal entre las columnas del Jacobiano de traslación.
    \item La dirección del movimiento local que se pierde (no puede ser ejecutada instantáneamente).
\end{enumerate}
\vspace{3cm}

\section*{Ejercicio 5 — Velocidad Cerca de la Singularidad[cite: 13]}
Considere el Jacobiano local $J = \begin{bmatrix} 1 & 0 \\ 0 & 0.05 \end{bmatrix}$ y una velocidad cartesiana deseada $\dot{\mathbf{x}}_d = \begin{bmatrix} 0 \\ 0.5 \end{bmatrix}$[cite: 13].
Calcule el vector de velocidades articulares requeridas $\dot{\mathbf{q}}$. Explique el efecto de "amplificación" de velocidad que ocurre físicamente.
\vspace{3cm}

\end{document}
